% =====================================================================
%  GMN3 §10.1--10.3  $A_1$ 理论的例子（Examples of $A_1$ theories）
%  原文：Lamination-Paper.tex 5360--5654 行；PDF p.86--92（印刷 p.85--91）
%  编号公式 (10.1)--(10.34) 共 34 条：9 个 equation + 5 个 align（25 行），
%  行结构与原文逐行一致，故式号一一对应；插图 14 幅，全部在本块内。
%  跨块引用（节号、式号、附录）暂写原文编号并加 % TODO，合稿时统一改 \ref/\eqref
% =====================================================================

\section{$A_1$ 理论的例子}\label{sec:Examples}

本节我们在 $A_1$ 理论的几个简单例子中导出带框 BPS 简并度。

我们还会评述这些理论的保护自旋特征标方面的结果。我们将借助下面这条简单的评注来断定这些 PSC，并给它冠以一个响亮的名号：\emph{提升原理}（Promotion Principle）。想法非常简单。若已知某个自旋特征标在 $y=1$ 处的特殊化等于 $1$，那么相应的表示必定是 $SU(2)$ 的平凡（即自旋为零的）表示的单份拷贝，于是其特征标恒等于 $1$（对一切 $y$ 皆然）。注意，倘若我们所知的仅仅是 PSC 落在 $SU(2)$ 的表示环中，这条原理显然就会失效；因此，只有在强正性猜想成立时，我们才能援引提升原理。在下文第 11 节% TODO 合稿改 \ref{sec:Quantum-Holonomy}
中，我们会稍微更系统地利用这条原理。

\subsection{$N=3$ Argyres--Douglas 理论}\label{subsec:ADN3-Laminations}

我们从 $N=3$ Argyres--Douglas 理论开始；第 4.3 节% TODO 合稿改 \ref{sec:formaln3}
已对它作过简要回顾。这一情形下用 Hitchin 系统给出的 $\CM$ 之实现，见文献 \cite{Gaiotto:2009hg} 的第 9.4.4 节。此时 $C = \IC\IP^1$，只有一个不规则奇点，从它发出 $5$ 条斯托克斯线；相应地，边界 $S^1$ 上有 $5$ 个标记点。这些标记点把边界分成 $5$ 个扇区，我们按沿圆周顺时针递增的顺序把它们编号为 $1$ 到 $5$。
\insfig[0.45\textwidth]{pentagon-bare}{曲线 $C = \IC\IP^1$，其中挖去了无穷远附近的一个圆盘，边界上有五个标记点。}

本理论中有一组特别的 $5$ 个叶层。先从图 \ref{fig:pentagon-l1} 所示的 $L_1$ 开始。
\insfig[0.45\textwidth]{pentagon-l1}{$N=3$ 理论中的叶层 $L_1$。}
$L_{i+1}$ 由 $L_i$ 把圆周逆时针旋转角度 $\frac{4 \pi}{5}$ 得到。把这一操作重复 $5$ 次便回到原来的叶层，故 $L_{i+5} = L_i$。

这些叶层 $L_i$ 中的每一个都给出一族由 $\zeta$ 参数化的线算子 $L_{i,\zeta}$。为计算维数 $\fro(\gamma; u, \zeta, y = 1)$，根据式 (6.14)，% TODO 合稿改 \eqref{eq:y1}
我们应当把诸 vev $\inprod{L_{i,\zeta}(u)}'$ 按坐标函数 $\tCY_\gamma$ 展开。这可以借助我们在附录 F% TODO 合稿改 \ref{app:traffic}
中描述的「traffic rule」算法来完成。我们先画出 WKB 三角剖分 $T_\WKB(\zeta, u)$。五边形的各种三角剖分在顶点的循环置换之下看起来都一样，所以即使不知道 $\zeta$ 与 $u$ 是什么，我们也知道 $T_\WKB(\zeta, u)$ 长什么样。我们按图 \ref{fig:pentagon} 所示标记各顶点。
\insfig[0.45\textwidth]{pentagon}{$N=3$ 理论中的一个三角剖分。五边形的五条边都是边界边。边界上的每个橙色记号现在都等同于一条从该不规则奇点发出的斯托克斯射线。蓝色点是三角剖分的顶点；我们把它们等同于反斯托克斯射线。}

在附录 F% TODO 合稿改 \ref{app:traffic}
中，我们对每条边 $E$ 定义了模空间 $\tCM$ 上的函数 $\tCY_E(\zeta)$。为方便起见，我们减轻记号负担：把函数 $\tCY_X$ 写作 $X$，把 $\tCY_Y$ 写作 $Y$。于是应用 traffic rule，我们很快得到诸 vev 按 $X$ 与 $Y$ 的展开：
\begin{align}
\inprod{L_1}' &= X, \label{eq:pentagon-vevs-1} \\
\inprod{L_2}' &= Y + XY, \\
\inprod{L_3}' &= \frac{1}{X} + \frac{Y}{X} + Y, \\
\inprod{L_4}' &= \frac{1}{X} + \frac{1}{XY}, \\
\inprod{L_5}' &= \frac{1}{Y}. \label{eq:pentagon-vevs-5}
\end{align}
由这些 vev，我们立刻读出带框 BPS 态的谱。例如，线算子 $L_1(\zeta)$ 在真空 $u$ 中只支持单个带框 BPS 态。这个态因而必定自旋为零。另一方面，$L_3(\zeta)$ 支持三个带框 BPS 态，三者电荷各不相同。同样，这些态的自旋都为零。

本理论中一般的叶层形如
\begin{equation}\label{eq:ADN3-gen-lam}
L_i^m L_{i+1}^n, \qquad m, n \in \IZ_{\ge 0}.
\end{equation}
（两个不相邻的 $L_i$ 之积不给出叶层，因为叶层的边不允许彼此相交。）相应的期望值为
\begin{equation}
\inprod{L_i}^m \inprod{L_{i+1}}^n, \qquad m, n \in \IZ_{\ge 0}.
\end{equation}
展开这样的期望值，一般会得到更有趣的 BPS 态谱，其中某些电荷扇区的重数大于 $1$。

计算这些算子的交换环很有意思。为此，只需给出五个生成元 $L_i$ 的环关系。检视式 \eqref{eq:pentagon-vevs-1}--\eqref{eq:pentagon-vevs-5} 可见 $L_{i+1} L_{i-1} = 1 + L_i$。（这也可以更直接地理解为「skein 关系」的推论：skein 关系把两个相交叶层之积，与一个把交点解消掉的叶层之和联系起来。）

注意，若作识别 $Y \to X_{-\gamma_1}$ 与 $X \to X_{\gamma_2}$，则表达式 \eqref{eq:pentagon-vevs-1}--\eqref{eq:pentagon-vevs-5} 与第 4.3 节% TODO 合稿改 \ref{sec:formaln3}
中腔 $c_1$ 里求值的形式线算子 $F_i$ 的相应表达式完全吻合。此外，式 \eqref{eq:ADN3-gen-lam} 恰好对应第 4.3 节% TODO 合稿改 \ref{sec:formaln3}
所定义的简单形式线算子 $F_{p,q}^{(i)}$。

\subsection{$N = 4$ Argyres--Douglas 理论}

现在考虑 $N=4$ Argyres--Douglas 理论。它用 Hitchin 系统的实现与我们刚讨论的 $N=3$ 情形非常相似：唯一的差别是边界 $S^1$ 上有 $6$ 个标记点而不是 $5$ 个。见图 \ref{fig:hexagon-bare}。
\insfig[0.45\textwidth]{hexagon-bare}{曲线 $C = \IC\IP^1$，其中挖去了无穷远附近的一个圆盘，边界上有六个标记点。}

与上一情形的一个重要差别来自 $6$ 是偶数：存在非平凡的叶层，可以被推入边界的一个任意小邻域。一个例子见图 \ref{fig:hexagon-mu}。
\insfig[0.45\textwidth]{hexagon-mu}{一个非平凡的叶层 $\mu$，它可以被推入边界的一个任意小邻域。}
$\inprod{\mu}$ 与 $\CM$ 上的位置无关：不过，它可以通过理论的非可归一化形变而改变。

现在按图 \ref{fig:hexagon-l1} 定义两个叶层 $L_1$ 与 $L_2$。
\insfig[0.62\textwidth]{hexagon-l1}{叶层 $L_1$ 与 $L_2$。}
再把「由 $L_i$ 把圆周顺时针旋转角度 $\frac{2 \pi}{6}$ 所得的叶层」定义为 $L_{i+2}$。于是序列 $L_i$ 以 $12$ 为周期。此外，$L_{i+6}$ 可以利用下述关系表成 $L_i$ 的式子：
\begin{align}
L_7 &= \mu^{-1} L_1, \\
L_8 &= L_2,
\end{align}
以及它们在操作 $(L_i \to L_{i+2}, \mu \to \mu^{-1})$ 作用下的像。任何叶层都具有形式
\begin{equation}
\mu^r L_i^m L_{i+1}^n, \quad r \in \IZ, \quad m, n \in \IZ_{\ge 0}
\end{equation}
（对某个 $i$）。

\insfig[0.45\textwidth]{hexagon}{$N=4$ Argyres--Douglas 理论中一个三角剖分的例子。}

现在考虑图 \ref{fig:hexagon} 所示的三角剖分；对某组 $(u, \zeta)$，它可能出现为 $T_\WKB(u, \zeta)$。把这些叶层按此三角剖分所附的坐标展开——仍然使用 traffic rule——得到
\begin{align}
\inprod{\mu}' &= XZ, \\
\inprod{L_1}' &= Z, \\
\inprod{L_2}' &= Y + YZ, \\
\inprod{L_3}' &= \frac{1}{Z} + \frac{Y}{Z} + Y, \\
\inprod{L_4}' &= \frac{(1+Y)(1 + X + XY + XYZ)}{YZ}, \\
\inprod{L_5}' &= \frac{1}{Y} + \frac{X}{Y} + X, \\
\inprod{L_6}' &= \frac{1}{XY} + \frac{1}{Y}.
\end{align}
利用这些显式公式，我们可以轻易读出诸 $L_i$ 之间的关系：我们有
\begin{align}
L_1 L_3 &= 1 + L_2, \\
L_2 L_4 &= (1+ L_3)(1 + L_9),\\
L_1 L_4 &= 1+ \mu + L_5 + L_9,
\end{align}
以及它们在操作 $(L_i \to L_{i+2}, \mu \to \mu^{-1})$ 作用下的像。

\subsection{$SU(2)$、$N_f = 0$}\label{subsec:SU2-Lams}

接下来处理纯 $SU(2)$ 规范理论。这是取 $C = \IC\IP^1$、并带两个不规则奇点的 $A_1$ 理论，$\phi_2$ 在每个奇点处都有三阶极点。拓扑上，正如我们已经提到的，我们应当在每个奇点周围挖去一个小圆盘，并把每条边界分支分成单独一段弧，即在每条边界分支上只标记一个点——见图 \ref{fig:nf0-bare}。

\insfig[0.62\textwidth]{nf0-bare}{曲线 $C = \IC\IP^1$，其中挖去了两个圆盘，每条边界分支上各有一个标记点。}

叶层定义如下。先取两个小半圆 $\wp_1$、$\wp_2$，分别越过两个边界圆周上的标记点。再取两条不相交的曲线 $\wp_3$、$\wp_4$，从一条边界圆周走到另一条，二者彼此相差一个单位的缠绕。这四条曲线构成一个「极大」集，意思是：我们无法再加入任何一条既不相交、又不与其中任何一条同痕的曲线。给它们配上整数权重 $k_i$，并服从限制
\begin{equation} \label{eq:nf0-cond}
2 k_1 + k_3 + k_4 = 0, \quad 2 k_2 + k_3 + k_4 = 0, \quad k_3 \ge 0, \quad k_4 \ge 0.
\end{equation}
这样我们就得到一族叶层，由式 \eqref{eq:nf0-cond} 的解集参数化。对每一对满足 $k_3 + k_4$ 为偶数的 $k_3$、$k_4$，我们得到一个解。见图 \ref{fig:nf0-maximal}。
\insfig[0.62\textwidth]{nf0-maximal}{一族叶层。}
这只是无穷多族此类叶层中的一族，因为我们可以任意选择 $\wp_3$、$\wp_4$ 的整体缠绕。

所有这些叶层都给出超对称线算子。它们究竟是哪些算子？把文献 \cite{Drukker:2009tz} 给出的对应关系应用于 $N_f = 4$ 的 $SU(2)$ 理论，再沿第 9.3 节% TODO 合稿改 \ref{subsec:Decouple-Flavor}
的思路把质量送到无穷大，就会发现这些叶层对应 Wilson--'t Hooft 算子。它们是标号为 $[(\frac{p}{2} H_\alpha , \frac{q}{2} \alpha)]$ 的算子，其中磁荷 $p = k_3 + k_4$，电荷 $q$ 等于总缠绕，即 $k_3$ 乘以 $\wp_3$ 的缠绕加上 $k_4$ 乘以 $\wp_4$ 的缠绕。\footnote{对一条从一条边界走到另一条边界的曲线而言，「总缠绕」的含义并不唯一：真正典范地定义了的，只是两条此类曲线之间的相对缠绕。这里的这种不唯一性，反映的是规范理论中定义 Wilson--'t Hooft 算子之电荷时类似的困难。}在单值变换 $\zeta \to e^{2 \pi i} \zeta$ 之下，$\wp_3$ 与 $\wp_4$ 都增加 $2$ 个单位的缠绕，于是 $q$ 平移 $2p$ 个单位。

还有少数叶层我们尚未描述：即，我们可以取一条把两个奇点分开的单个闭曲线，它承载 $SL(2,\IC)$ 的自旋 $\frac{q}{2}$ 表示。这个叶层对应一个纯粹的 Wilson 圈算子，其电荷为 $q$，亦即处于最高权为 $\frac{q}{2} \alpha$ 的表示之中。

现在可以来考虑诸 vev 的红外展开了。照例，我们先画出 WKB 三角剖分 $T_\WKB(u, \zeta)$。此时组合结构非常简单，任何这样的三角剖分都必定形如图 \ref{fig:nf0}。（特别地，无论 $u$ 在弱耦合区还是强耦合区，都是如此！）有两条内部边，我们分别标为 $X$ 与 $Y$，二者相差一个单位的缠绕。
%
\insfig[0.62\textwidth]{nf0}{$SU(2)$、$N_f = 0$ 理论中的一个三角剖分。}
设 $L_0$ 与 $L_1$ 是图 \ref{fig:nf0-l0} 与图 \ref{fig:nf0-l1} 所示的叶层。更一般地，我们可以从 $L_n$ 出发、绕两个边界 $S^1$ 之一作 $k$ 次 Dehn twist，以此来定义 $L_{n+2k}$。这样就得到对一切 $n \in \IZ$ 的叶层 $L_n$。例如 $L_2$ 见图 \ref{fig:nf0-l2}。我们有
\begin{equation}
L_{2n+1}^2 = L_{2n} L_{2n+2}.
\end{equation}
（要验证这一点，只需注意：借助 Dehn twist 可以化归到 $n=0$ 的情形，而该情形可由图 \ref{fig:nf0-l0}--\ref{fig:nf0-l2} 看出。）
\insfig[0.62\textwidth]{nf0-l0}{$SU(2)$、$N_f = 0$ 理论中的叶层 $L_0$。}
\insfig[0.62\textwidth]{nf0-l1}{$SU(2)$、$N_f = 0$ 理论中的叶层 $L_1$。}
\insfig[0.62\textwidth]{nf0-l2}{$SU(2)$、$N_f = 0$ 理论中的叶层 $L_2$。}

最一般的叶层形如 $L_n^a L_{n+1}^b$，或者是一条承载 $SL(2,\IC)$ 某个表示的闭合回路。我们用 $L_*$ 记由单个闭合回路组成的叶层，该回路承载 $SL(2,\IC)$ 的基础表示。见图 \ref{fig:nf0-lstar}。
\insfig[0.62\textwidth]{nf0-lstar}{$SU(2)$、$N_f = 0$ 理论中的叶层 $L_*$。注意本图画在 $\IC\IP^1$ 上，因此，尽管我们把圆画成绕过左侧不规则奇点的样子，它其实与绕过右侧奇点的圆同痕。这里我们对记号稍有滥用：用 $+1$ 代表 $SL(2,\IC)$ 的基础表示。}

现在我们要计算与这些叶层相联系的线算子的期望值 $\inprod{L}'$，并把它们按 $\tCM$ 上的函数 $\tCY_X$ 与 $\tCY_Y$ 展开。照例，下文把这两个函数简记为 $X$ 与 $Y$。对任何具体叶层，结果都可以由 traffic rule 直接得到。然而事实上，答案还有一个漂亮的统一公式，它可以借更直接地利用式 (9.11)% TODO 合稿改 \eqref{eq:open-vev}
来证明。我们把该推导放在附录 G% TODO 合稿改 \ref{app:chebyshev}
中，这里只报告答案。

引入由 $U_n(\cos\theta) := \frac{\sin(n+1)\theta}{\sin\theta}$ 与 $T_n(\cos\theta) := \cos(n\theta)$ 定义的切比雪夫多项式（Tchebyshev）。令
\begin{equation}
 \alpha := \frac{1}{2 \sqrt{XY}} \left( XY + Y + 1 \right),
\end{equation}
则公式为
\begin{equation}\label{eq:chebyshev-form}
\inprod{L_{2n}}' = Y^{-1} \left[T_n(\alpha) + \frac{XY- Y - 1}{2 \sqrt{X Y}} U_{n-1}(\alpha) \right]^2.
\end{equation}
特别地，所有期望值 $\inprod{L_{2n}}'$ 都可以表为 $X$ 与 $Y$ 的、系数为正整数的洛朗多项式，正如我们所预期的。例如，
\begin{align}
\inprod{L_{-4}}' & = X^{-2} Y^{-3} (1+ 2 Y + Y^2 + X Y^2)^2 \label{eq:chebyshev-expls} \\
&= \frac{1}{Y^3 X^2}+\frac{4}{Y^2 X^2}+\frac{Y}{X^2}+\frac{6}{Y X^2}+\frac{2 Y}{X}+\frac{2}{Y X}+Y+\frac{4}{X^2}+\frac{4}{X}, \\
\inprod{L_{-2}}' &= \frac{1}{X Y^2} (1 + Y)^2=\frac{1}{Y^2 X}+\frac{2}{Y X}+\frac{1}{X}, \\
\inprod{L_{0}}' &= \frac{1}{Y}, \\
\inprod{L_{2}}' &= X, \\
\inprod{L_{4}}' &= Y(1+X)^2 = Y X^2+2 Y X+Y, \\
\inprod{L_{6}}' & = X^{-1} (1+ (1+X)^2 Y)^2 \\
&= Y^2 X^3+4 Y^2 X^2+6 Y^2 X+\frac{Y^2}{X}+4 Y^2+2 Y X+\frac{2 Y}{X}+4 Y+\frac{1}{X}, \label{eq:chebyshev-expll}
\end{align}
如此等等。我们还有
%
\begin{equation} \label{eq:wilson-vev-nf0}
\inprod{L_*}' = \frac{1}{\sqrt{XY}} + \sqrt{\frac{Y}{X}} + \sqrt{XY}.
\end{equation}
根据式 (6.14)，% TODO 合稿改 \eqref{eq:y1}
这些展开捕获了带框 BPS 简并度。例如，在式 \eqref{eq:wilson-vev-nf0} 中我们发现，插入了基础表示紫外 Wilson 圈的希尔伯特空间里有三个带框 BPS 态。其中 $(XY)^{\pm \half}$ 两项对应带电荷的态，电荷为 $\pm 1$，这与我们朴素预期的相符。更令人惊讶的是，还有一项额外贡献 $\sqrt{X/Y}$，它对应一个不带电荷的态，其磁荷由边 $X$ 与边 $Y$ 之缠绕的差给出。从弱耦合规范理论的计算中重现这一结果，将会很有意思。此外，展开式 \eqref{eq:chebyshev-expls}--\eqref{eq:chebyshev-expll} 蕴涵了一幅束缚在紫外 't Hooft--Wilson 圈上的带框 BPS 态的复杂图景；从弱耦合的角度发现它，应当也很有意思。

线算子之间的交换环关系也值得考察。例如，式 (G.4)% TODO 合稿改 \eqref{eq:Trig-Ln}
的一个简单推论是
%
\begin{equation} L_{2k} L_* = L_{2k-1} + L_{2k+1}. \end{equation}
%
这在物理上非常合理：若把一个 Wilson 圈与一个 't Hooft--Wilson 圈放到一起，后者把 $SU(2)$ 破缺为 $U(1)$，于是 Wilson 圈分解为两个电荷相反的 $U(1)$ 贡献之和。这样我们就得到两个 't Hooft--Wilson 圈，其电荷平移了一个单位。我们还有关系 $L_{2k} L_{2k+4} = (1+L_{2k+2})^2$、$L_{2k} L_{2k+6} = 1+(1+ L_{2k+2})(1 + L_{2k+4}) + L_*^2$、$L_{2k+1} L_{2k+3} = L_{2k+2} (1+L_{2k+2})$，等等。这些关系很难在规范理论中直接理解。

最后，与第 4.4 节% TODO 合稿改 \ref{subsec:formal-SU2}
中的形式线算子作比较很有意思。我们辨认出：在识别 $X=[0,2]$ 与 $Y=[-2,0]$ 之下，可以把 $L_{2n}$ 等同于 $G_{2n-1}$，从而把 $\sqrt{L_{2n}}$ 等同于在 $y=1$ 处求值的 $\hat V_{n-1}$。上述诸环关系就成为式 (4.50)--(4.54)% TODO 合稿改 \eqref{eq:SU2-formring-1}--\eqref{eq:SU2-formring-5}
在 $y=1$ 处的那些关系。这一比较与第 8.3 节% TODO 合稿改 \ref{subsec:SU2-Functions}
末尾的讨论相容，并且进一步表明：我们可以用形式线算子显式地算出这些理论中的保护自旋特征标。注意，一般说来，并不能简单地把式 \eqref{eq:chebyshev-expls} 中的整数 $n$ 提升为 $[n]$。
